\chapter{The Twelve-Key Pattern Hiding Inside Eighty-Eight Keys}
\label{ch:keyboard-geometry}
\chapterline

Stand in front of a full-sized grand or upright piano for the first time and your eyes meet eighty-eight black and white keys stretching nearly four feet wide. Looking at all eighty-eight keys at once feels like staring at a yardstick with eighty-eight unnumbered tick marks. Where do your hands even belong?

Take two steps back and squint slightly at the raised black keys. A calm, repeating rhythm jumps out right away:
\[
  \text{two black keys, gap, three black keys, gap, two black keys, gap, three black keys.}
\]
An eighty-eight-key piano is not eighty-eight different notes. It is \textbf{one twelve-key tile} --- seven white keys and five raised black keys --- copied seven full times side by side, bookended by three extra low keys on the far left and one final high key on the far right ($3 + 7 \times 12 + 1 = 88$). Understand that single twelve-key tile once, and you understand the entire length of the instrument.

\section{Instant Landmarks Without Stickers}

Many beginners paste paper letter stickers onto their white keys. Within a week, their eyes are glued to the stickers rather than listening to the sound, and sitting down at any unlabelled piano feels like walking into a dark room. You never need stickers because the raised black keys hand you two permanent physical landmarks by sight and by touch.

Look at the two black-key clusters inside each twelve-key tile:
\begin{itemize}
  \item \textbf{The Chopsticks (group of 2 black keys):} Slide your fingertip onto the white key sitting immediately to the \textbf{left} of any pair of two black keys. That note is always \textbf{$C$}.
  \item \textbf{The Fork (group of 3 black keys):} Slide your fingertip onto the white key sitting immediately to the \textbf{left} of any group of three black keys. That note is always \textbf{$F$}.
\end{itemize}

Once your eye spots $C$ and $F$, the remaining five white keys fill in alphabetically from left to right using only the first seven letters of the alphabet ($A$ through $G$):
\[
  C \;\to\; D \;\to\; E \;\to\; F \;\to\; G \;\to\; A \;\to\; B \;\to\; C.
\]
After $G$, the musical alphabet wraps straight back to $A$ (never $H$). Notice where each white key sits in Figure~\ref{fig:two-octave-blueprint}. The note $D$ sits snugly between the two chopsticks, while $E$ sits right after them. Similarly, $G$ and $A$ sit inside the three-pronged fork, and $B$ sits right after the fork just before the next $C$.

\begin{figure}[htbp]
\centering
\begin{tikzpicture}
  % Two-octave blueprint: 14 white keys (C4 to B5), width w = 0.88cm -> 12.32cm total width!
  \def\wk{0.88}
  \def\hk{2.80}

  % Highlight C4 (k=0), F4 (k=3), C5 (k=7), F5 (k=10)
  \foreach \k in {0,...,13} {
    \filldraw[fill=white, draw=black!65, line width=0.6pt] ({\k*\wk},0) rectangle ({(\k+1)*\wk},\hk);
  }
  \foreach \k in {0,7} {
    \filldraw[fill=sourceblue!15, draw=black!65, line width=0.6pt] ({\k*\wk},0) rectangle ({(\k+1)*\wk},\hk);
  }
  \foreach \k in {3,10} {
    \filldraw[fill=targetgreen!16, draw=black!65, line width=0.6pt] ({\k*\wk},0) rectangle ({(\k+1)*\wk},\hk);
  }

  % Black keys: octave 1 (0,1 and 3,4,5), octave 2 (7,8 and 10,11,12)
  \foreach \k in {0,1,3,4,5,7,8,10,11,12} {
    \filldraw[fill=black!85, draw=black, line width=0.5pt] ({(\k+1)*\wk-0.25},1.08) rectangle ({(\k+1)*\wk+0.25},\hk);
  }

  % White-key note names (no duplicate C/F nodes!)
  \foreach \k/\nm/\col in {
    0/C4/sourceblue!90!black,
    1/D/black!85,
    2/E/black!85,
    3/F4/targetgreen!80!black,
    4/G/black!85,
    5/A/black!85,
    6/B/black!85,
    7/C5/sourceblue!90!black,
    8/D/black!85,
    9/E/black!85,
    10/F5/targetgreen!80!black,
    11/G/black!85,
    12/A/black!85,
    13/B/black!85%
  } {
    \node[font=\footnotesize\bfseries, text=\col] at ({(\k+0.5)*\wk},0.34) {\nm};
  }

  % Semitone ruler 0..11 cleanly ABOVE Octave 1 (zero vertical seams behind digits!)
  \foreach \xpos/\idx in {
    0.5/0, 1.0/1, 1.5/2, 2.0/3, 2.5/4,
    3.5/5, 4.0/6, 4.5/7, 5.0/8, 5.5/9, 6.0/10, 6.5/11%
  } {
    \draw[gray!60, line width=0.45pt] ({\xpos*\wk},2.82) -- ({\xpos*\wk},2.94);
    \node[font=\scriptsize\bfseries, text=structurepurple!90!black] at ({\xpos*\wk},3.14) {\idx};
  }

  % Brackets above black-key groups (raised above semitone ruler on Octave 1, crisp on Octave 2)
  \draw[decorate, decoration={brace, amplitude=5pt}, thick, sourceblue]
    ({1*\wk-0.26},3.40) -- ({2*\wk+0.26},3.40)
    node[midway, above=5pt, font=\footnotesize\bfseries, text=sourceblue] {2: Chopsticks};
  \draw[decorate, decoration={brace, amplitude=5pt}, thick, targetgreen!80!black]
    ({4*\wk-0.26},3.40) -- ({6*\wk+0.26},3.40)
    node[midway, above=5pt, font=\footnotesize\bfseries, text=targetgreen!80!black] {3: Fork};
  \draw[decorate, decoration={brace, amplitude=5pt}, thick, sourceblue]
    ({8*\wk-0.26},2.95) -- ({9*\wk+0.26},2.95)
    node[midway, above=5pt, font=\footnotesize\bfseries, text=sourceblue] {2: Chopsticks};
  \draw[decorate, decoration={brace, amplitude=5pt}, thick, targetgreen!80!black]
    ({11*\wk-0.26},2.95) -- ({13*\wk+0.26},2.95)
    node[midway, above=5pt, font=\footnotesize\bfseries, text=targetgreen!80!black] {3: Fork};

  % Below-keyboard annotations: 12-key tile bracket and natural half-step seams
  \draw[decorate, decoration={brace, mirror, amplitude=6pt}, thick, structurepurple!85!black]
    (0.04,-0.16) -- ({7*\wk-0.04},-0.16)
    node[midway, below=7pt, font=\footnotesize\bfseries, text=structurepurple!85!black] {One 12-Key Tile (Semitones $0$ through $11$)};

  % Natural half-step gap callouts on Octave 2 (E5-F5 at 10*wk, B4-C5 at 7*wk)
  \node[draw=quotientorange!85!black, fill=quotientorange!8, rounded corners=3pt, inner xsep=6pt, inner ysep=4pt, font=\scriptsize, align=center] (seam) at ({10.5*\wk},-0.72)
    {Natural white-to-white seams:\\no black key between $B$--$C$ or $E$--$F$};
  \draw[->, thick, >=Latex, quotientorange!85!black, shorten >=2pt] (seam.north west) -- ({7*\wk},-0.04);
  \draw[->, thick, >=Latex, quotientorange!85!black, shorten >=2pt] (seam.north) -- ({10*\wk},-0.04);
\end{tikzpicture}
\caption{Two repeating twelve-key tiles from Middle $C$ ($C4$) up to $B5$. Every $C$ sits immediately left of the two black chopsticks; every $F$ sits immediately left of the three-pronged fork. Semitone indices $0$ through $11$ count all twelve keys in order inside each tile.}
\label{fig:two-octave-blueprint}
\end{figure}

Musicians attach a small octave number to tell one $C$ from the next. The $C$ closest to the middle of the piano bench (usually right underneath the brand name on the fallboard) is called \textbf{Middle $C$}, written $C4$. Moving right to the next $C$ gives $C5$; moving left gives $C3$, $C2$, and $C1$, where notes rumble deep in the bass.

We also number the five fingers on \emph{both} hands outward from the thumbs:
\begin{center}
\small
\begin{tabular}{@{}ccccc@{}}
\toprule
\textbf{1} & \textbf{2} & \textbf{3} & \textbf{4} & \textbf{5} \\
\midrule
Thumb & Index finger & Middle finger & Ring finger & Little finger (pinky) \\
\bottomrule
\end{tabular}
\end{center}
Because your hands are mirror images of each other, finger~$1$ is always the thumb on \emph{both} hands, and finger~$5$ is always the pinky.

\begin{labbox}
\textbf{Blindfolded landmark drill (60 seconds):} Close your eyes in front of a keyboard. Lightly brush the raised black keys with fingers $2$, $3$, and $4$ until you feel a cluster of \textbf{two} black keys. Slide your thumb down onto the smooth white key immediately to their left, play it, and open your eyes to verify that you landed on $C$. Repeat three times across low, middle, and high octaves, then try the same touch test for the group of \textbf{three} black keys to land on $F$.
\end{labbox}

\section{Half Steps and Whole Steps: Look at the Back of the Keys}

Look along the \emph{front} edge of the piano and you see a solid row of white keys, which makes it tempting to think that stepping from one white key to the next is always the same musical distance. Now shift your gaze two inches deeper, toward the \textbf{back edge} of the keys near the wooden fallboard. Back there, narrow white tails and raised black keys line up in one even row of twelve strips per octave.

That back rail shows the true unit of distance on the piano:
\begin{itemize}
  \item A \textbf{half step} (also called a \textbf{semitone}) moves \textbf{one physical key} to the left or right along the back rail, with \emph{zero} keys skipped in between. Most half steps travel from a white key up to its neighbor black key (such as $C \to C\sharp$) or from a black key to a white key. However, where there is \textbf{no black key} in the gap --- right across $E$--$F$ and $B$--$C$ --- stepping directly from white key to white key is \emph{also} one half step.
  \item A \textbf{whole step} (also called a \textbf{whole tone}) equals \textbf{two half steps}: you skip over \textbf{one physical key} (black or white) along the back rail. Going from $C$ to $D$ skips the black key $C\sharp$, so $C \to D$ is a whole step. Going from $E$ up a whole step must skip the adjacent white key $F$ and land on the black key $F\sharp$, as Figure~\ref{fig:half-whole-steps} shows.
\end{itemize}

\begin{whybox}
Why did instrument makers leave out a black key between $E$--$F$ and $B$--$C$ instead of alternating white and black keys forever? Imagine an eighty-eight-key keyboard of strictly alternating white-black-white-black keys. Every key would look and feel identical to every other key; you could never spot $C$ or $F$ without counting from the far end of the piano. Grouping the twelve semitones into seven white keys (an ancient seven-note vocal scale) and five raised black keys gives your eyes and fingers an unmistakable map.
\end{whybox}

\begin{figure}[htbp]
\centering
\begin{tikzpicture}
  % Left keyboard: Half Steps (+1 physical key)
  \begin{scope}[shift={(0,0)}]
    \def\wk{0.76}
    \def\hk{2.45}
    \node[anchor=south, font=\footnotesize\bfseries, text=sourceblue] at ({4*\wk},3.35)
      {Panel (a): Half Step = $+1$ Key (Skip $0$)};
    \foreach \k in {0,...,7} {
      \filldraw[fill=white, draw=black!65, line width=0.6pt] ({\k*\wk},0) rectangle ({(\k+1)*\wk},\hk);
    }
    % Highlight E (k=2) and F (k=3) for white-to-white half step, and C (k=0)
    \filldraw[fill=sourceblue!14, draw=black!65, line width=0.6pt] ({0*\wk},0) rectangle ({1*\wk},\hk);
    \filldraw[fill=targetgreen!16, draw=black!65, line width=0.6pt] ({2*\wk},0) rectangle ({3*\wk},\hk);
    \filldraw[fill=targetgreen!16, draw=black!65, line width=0.6pt] ({3*\wk},0) rectangle ({4*\wk},\hk);

    \foreach \k in {0,1,3,4,5} {
      \filldraw[fill=black!85, draw=black, line width=0.5pt] ({(\k+1)*\wk-0.22},0.98) rectangle ({(\k+1)*\wk+0.22},\hk);
    }
    % Highlight C# (k=0 black key)
    \filldraw[fill=sourceblue!90, draw=black, line width=0.6pt] ({1*\wk-0.22},0.98) rectangle ({1*\wk+0.22},\hk);

    % Dots
    \fill[sourceblue] ({0.5*\wk},0.46) circle (0.11);
    \fill[white] ({1.0*\wk},1.32) circle (0.08);
    \fill[targetgreen!80!black] ({2.5*\wk},0.46) circle (0.11);
    \fill[targetgreen!80!black] ({3.5*\wk},0.46) circle (0.11);

    % Labels below white keys (clear of strokes)
    \node[font=\footnotesize\bfseries, text=sourceblue] at ({0.5*\wk},-0.30) {$C$};
    \node[font=\footnotesize\bfseries, text=targetgreen!80!black] at ({2.5*\wk},-0.30) {$E$};
    \node[font=\footnotesize\bfseries, text=targetgreen!80!black] at ({3.5*\wk},-0.30) {$F$};

    % Curved arrows ABOVE keyboard for C -> C#, BELOW keyboard for E -> F
    \draw[->, very thick, >=Latex, sourceblue, shorten <=2pt, shorten >=2pt]
      ({0.42*\wk},2.50) to[out=55,in=125] node[above=2pt, font=\scriptsize\bfseries, text=sourceblue] {$+1$ ($C \to C\sharp$)} ({1.0*\wk},2.50);
    \draw[->, very thick, >=Latex, targetgreen!80!black, shorten <=2pt, shorten >=2pt]
      ({2.5*\wk},-0.52) to[out=-45,in=-135] node[below=2pt, font=\scriptsize\bfseries, text=targetgreen!80!black] {$+1$ ($E \to F$: no black key)} ({3.5*\wk},-0.52);
  \end{scope}

  % Right keyboard: Whole Steps (+2 physical keys, skip 1)
  \begin{scope}[shift={(7.15,0)}]
    \def\wk{0.76}
    \def\hk{2.45}
    \node[anchor=south, font=\footnotesize\bfseries, text=quotientorange!85!black] at ({4*\wk},3.35)
      {Panel (b): Whole Step = $+2$ Keys (Skip $1$)};
    \foreach \k in {0,...,7} {
      \filldraw[fill=white, draw=black!65, line width=0.6pt] ({\k*\wk},0) rectangle ({(\k+1)*\wk},\hk);
    }
    % Highlight C (k=0), D (k=1), E (k=2)
    \filldraw[fill=sourceblue!14, draw=black!65, line width=0.6pt] ({0*\wk},0) rectangle ({1*\wk},\hk);
    \filldraw[fill=sourceblue!14, draw=black!65, line width=0.6pt] ({1*\wk},0) rectangle ({2*\wk},\hk);
    \filldraw[fill=quotientorange!16, draw=black!65, line width=0.6pt] ({2*\wk},0) rectangle ({3*\wk},\hk);

    \foreach \k in {0,1,3,4,5} {
      \filldraw[fill=black!85, draw=black, line width=0.5pt] ({(\k+1)*\wk-0.22},0.98) rectangle ({(\k+1)*\wk+0.22},\hk);
    }
    % Highlight F# (k=3 black key)
    \filldraw[fill=quotientorange!90, draw=black, line width=0.6pt] ({4*\wk-0.22},0.98) rectangle ({4*\wk+0.22},\hk);

    % Dots
    \fill[sourceblue] ({0.5*\wk},0.46) circle (0.11);
    \fill[sourceblue] ({1.5*\wk},0.46) circle (0.11);
    \fill[quotientorange!85!black] ({2.5*\wk},0.46) circle (0.11);
    \fill[white] ({4.0*\wk},1.32) circle (0.08);
    \fill[gray!65] ({3.5*\wk},0.46) circle (0.08);

    % Labels below
    \node[font=\footnotesize\bfseries, text=sourceblue] at ({0.5*\wk},-0.30) {$C$};
    \node[font=\footnotesize\bfseries, text=sourceblue] at ({1.5*\wk},-0.30) {$D$};
    \node[font=\footnotesize\bfseries, text=quotientorange!85!black] at ({2.5*\wk},-0.30) {$E$};
    \node[font=\scriptsize, text=gray!75!black] at ({3.62*\wk},-0.30) {(skip $F$)};

    % Arrows
    \draw[->, very thick, >=Latex, sourceblue, shorten <=2pt, shorten >=2pt]
      ({0.5*\wk},-0.52) to[out=-45,in=-135] node[below=2pt, font=\scriptsize\bfseries, text=sourceblue] {$+2$ (skips $C\sharp$)} ({1.5*\wk},-0.52);
    \draw[->, very thick, >=Latex, quotientorange!85!black, shorten <=2pt, shorten >=2pt]
      ({2.5*\wk},2.50) to[out=45,in=135] node[above=2pt, font=\scriptsize\bfseries, text=quotientorange!85!black] {$+2$ ($E \to F\sharp$, skips $F$)} ({4.0*\wk},2.50);
  \end{scope}
\end{tikzpicture}
\caption{Counting along the back edge of the keys. Left: Every half step moves $+1$ physical key with nothing skipped --- whether white-to-black ($C \to C\sharp$) or white-to-white across a missing black key ($E \to F$). Right: Every whole step moves $+2$ half steps, skipping one key in between ($C \to D$ skips $C\sharp$; $E \to F\sharp$ skips the white key $F$).}
\label{fig:half-whole-steps}
\end{figure}

\begin{trapbox}
The single most common beginner mistake at the piano is assuming that ``whole step means white key to white key.'' Test $E \to F$ and $B \to C$ along the back edge: zero keys sit between them, so both are \textbf{half steps}. Whenever you count a whole step up from $E$ or $B$, you \emph{must} skip the neighboring white key and land on the black key $F\sharp$ or $C\sharp$.
\end{trapbox}

\section{Sharps ($\sharp$) and Flats ($\flat$) Are Directions, Not Separate Notes}

Beginners sometimes treat ``sharp notes'' and ``flat notes'' as mysterious extra pitches reserved for advanced players. In truth, $\sharp$ and $\flat$ are just one-step directional arrows:
\begin{itemize}
  \item \textbf{Sharp ($\sharp$)} means \textbf{step $1$ half step to the right} (one physical key higher in pitch).
  \item \textbf{Flat ($\flat$)} means \textbf{step $1$ half step to the left} (one physical key lower in pitch).
\end{itemize}

Look at the black key sitting between $F$ and $G$. Start on $F$ and step one half step right: you arrive at $F\sharp$. Now start on $G$ and step one half step left: you arrive at $G\flat$ --- the \emph{exact same physical key}. Musicians call two names for the same physical key \textbf{enharmonic equivalents}. Calling one street corner ``just east of 5th Avenue'' ($F\sharp$) or ``just west of 6th Avenue'' ($G\flat$) describes the exact same patch of pavement from two neighboring landmarks. In Chapter~2 we will see how choosing between $F\sharp$ and $G\flat$ keeps scales clean so every letter from $A$ to $G$ appears once.

\begin{recoverybox}
\textbf{Pause and test the rule:} Because $\sharp$ simply means ``$+1$ physical key right,'' what key on the piano is $E\sharp$? Put your finger on $E$ and move one physical key to the right: you land directly on the white key $F$. Likewise, $C\flat$ is one key left of $C$, which is the white key $B$. Sharps and flats do not mean ``black keys only'' --- they mean $+1$ key and $-1$ key from wherever you started.
\end{recoverybox}

\section{Why Music Sounds Good: Wave Ratios and the Twelve-Step Clock}

Press $C$ and $G$ together on a piano and the two notes blend into a calm, steady ring. Press $C$ and $C\sharp$ together right beside each other and the air buzzes with rough, restless friction. Why do certain key distances blend smoothly while others clash?

When a piano hammer strikes a steel string, the string vibrates back and forth hundreds of times per second, pushing waves of air toward your eardrum. For example, the $A$ above Middle~$C$ ($A4$) vibrates at $440$~cycles per second ($440\text{ Hz}$). Two notes sound settled (\textbf{consonant}) when their vibration rates lock into a \textbf{ratio of small whole numbers}, so their airwaves line up in a short, clean repeating cycle:
\begin{enumerate}
  \item \textbf{The Octave ($2:1$ ratio, $+12$ half steps):} Double $440\text{ Hz}$ to $880\text{ Hz}$ ($A5$), or halve it to $220\text{ Hz}$ ($A3$). Two crests of the higher wave align with every single crest of the lower wave. Your ear hears the two waves fuse so cleanly that we give both keys the \emph{same letter name} ($A$).
  \item \textbf{The Perfect Fifth ($3:2 = 1.50$ ratio, $+7$ half steps):} Multiply a string's frequency by $\frac{3}{2}$ (for example, $220\text{ Hz} \to 330\text{ Hz}$, which is $A \to E$, or $C \to G$). Every three waves of the upper note lock cleanly onto two waves of the lower note, producing an open, grounded ring.
  \item \textbf{The Major Third ($5:4 = 1.25$ ratio, $+4$ half steps):} Multiply frequency by $\frac{5}{4}$ (such as $C \to E$). Five upper waves align with four lower waves, adding warmth and brightness.
\end{enumerate}

\begin{historybox}
Ancient Greek, Chinese, and Arabic string makers tuned harps and lutes by stacking pure $3:2$ fifths one after another. Twelve pure fifths almost return to seven pure octaves, missing by about $1.36\%$ because $(3/2)^{12} \approx 129.75$ whereas $2^7 = 128$ (no power of $\frac{3}{2}$ can ever equal a power of $2$). For centuries, hiding that tiny leftover gap meant some keys on a harpsichord sounded sweet while distant keys howled out of tune. Dividing each $2:1$ octave into twelve \emph{equal} multiplicative steps --- \textbf{twelve-tone equal temperament}, celebrated by J.~S.~Bach's \emph{Well-Tempered Clavier} (1722) and standard on modern pianos --- spreads that tiny gap invisibly across all twelve keys so every key plays cleanly.
\end{historybox}

\begin{formalbox}
In \textbf{twelve-tone equal temperament}, stepping up one half step multiplies frequency by the constant factor
\[
  r = 2^{1/12} \approx 1.059463 \qquad (\text{about } +5.95\% \text{ per key}).
\]
Stepping up $k$ half steps multiplies frequency by $r^k = 2^{k/12}$:
\begin{itemize}
  \item $k = 12$ half steps (one octave): $2^{12/12} = 2.0000$ (exact $2:1$ octave).
  \item $k = 7$ half steps (perfect fifth, e.g.\ $C \to G$): $2^{7/12} \approx 1.4983$ (within $0.11\%$ of $\frac{3}{2} = 1.5000$).
  \item $k = 4$ half steps (major third, e.g.\ $C \to E$): $2^{4/12} \approx 1.2599$ (close to $\frac{5}{4} = 1.2500$).
\end{itemize}
Because every $+12$ steps doubles frequency and repeats the note name, the twelve pitch classes form the modular clock $\Z_{12} = \set{0, 1, 2, \dots, 11}$ shown in Figure~\ref{fig:ratios-and-clock}. Shifting any melody or chord up by $s$ half steps maps every key $x$ to $(x + s) \bmod 12$, preserving every relative ratio $2^{((x+s)-(y+s))/12} = 2^{(x-y)/12}$.
\end{formalbox}

\begin{figure}[htbp]
\centering
\begin{tikzpicture}
  % Left panel: Simple Vibrating String Ratios
  \begin{scope}[shift={(0,0)}]
    \node[anchor=south, font=\small\bfseries, text=sourceblue] at (3.20,4.85)
      {Panel (a): Simple Wave Ratios};

    % Fundamental 1:1 at y=3.75
    \node[anchor=east, font=\footnotesize\bfseries, text=sourceblue] at (0.85,3.75) {$1:1$};
    \draw[gray!55, dashed] (1.00,3.75) -- (5.80,3.75);
    \draw[sourceblue, very thick, domain=1.00:5.80, samples=80] plot (\x, {3.75 + 0.42*sin(180*(\x-1.00)/4.80)});
    \fill[black!80] (1.00,3.75) circle (0.06);
    \fill[black!80] (5.80,3.75) circle (0.06);
    \node[font=\scriptsize, text=gray!80!black] at (3.40,3.05) {Home string ($C4$, $+0$ keys)};

    % Octave 2:1 at y=2.15
    \node[anchor=east, font=\footnotesize\bfseries, text=targetgreen!80!black] at (0.85,2.15) {$2:1$};
    \draw[gray!55, dashed] (1.00,2.15) -- (5.80,2.15);
    \draw[targetgreen!80!black, very thick, domain=1.00:5.80, samples=100] plot (\x, {2.15 + 0.40*sin(360*(\x-1.00)/4.80)});
    \fill[black!80] (1.00,2.15) circle (0.06);
    \fill[black!80] (3.40,2.15) circle (0.06);
    \fill[black!80] (5.80,2.15) circle (0.06);
    \node[font=\scriptsize, text=gray!80!black] at (3.40,1.45) {Octave ($C5$, $+12$ keys: $2^{12/12}=2$)};

    % Fifth 3:2 at y=0.55
    \node[anchor=east, font=\footnotesize\bfseries, text=quotientorange!85!black] at (0.85,0.55) {$3:2$};
    \draw[gray!55, dashed] (1.00,0.55) -- (5.80,0.55);
    \draw[quotientorange!85!black, very thick, domain=1.00:5.80, samples=120] plot (\x, {0.55 + 0.38*sin(540*(\x-1.00)/4.80)});
    \fill[black!80] (1.00,0.55) circle (0.06);
    \fill[black!80] (2.60,0.55) circle (0.06);
    \fill[black!80] (4.20,0.55) circle (0.06);
    \fill[black!80] (5.80,0.55) circle (0.06);
    \node[font=\scriptsize, text=gray!80!black] at (3.40,-0.15) {Fifth ($G4$, $+7$ keys: $2^{7/12}\approx 1.498$)};
  \end{scope}

  % Right panel: 12-Note Chromatic Clock Z_12 centered at (10.55, 2.30)
  \begin{scope}[shift={(10.55,2.30)}]
    \node[anchor=south, font=\small\bfseries, text=structurepurple!85!black] at (0,2.55)
      {Panel (b): The Chromatic Clock ($\Z_{12}$)};

    % Inner circle of radius 1.42cm; tick marks outward to 1.56cm; labels safely outside at radius 2.02cm!
    \draw[thick, gray!55] (0,0) circle (1.38);

    % Highlight C(0) - E(4) - G(7) triangle inside the circle (vertices at radius 1.24cm, well inside circle!)
    \filldraw[fill=sourceblue!14, draw=sourceblue, very thick]
      ({90-0*30}:1.24) -- ({90-4*30}:1.24) -- ({90-7*30}:1.24) -- cycle;

    \foreach \idx/\lbl in {
      0/{$0$: $C$},
      1/{$1$: $C\sharp$},
      2/{$2$: $D$},
      3/{$3$: $E\flat$},
      4/{$4$: $E$},
      5/{$5$: $F$},
      6/{$6$: $F\sharp$},
      7/{$7$: $G$},
      8/{$8$: $A\flat$},
      9/{$9$: $A$},
      10/{$10$: $B\flat$},
      11/{$11$: $B$}
    } {
      \pgfmathsetmacro{\ang}{90 - \idx*30}
      \draw[thick, gray!70] (\ang:1.30) -- (\ang:1.46);
      \node[font=\scriptsize\bfseries, text=black!85] at (\ang:1.96) {\lbl};
    }
    % Emphasize vertices 0, 4, 7 with clean dots on inner ring
    \fill[sourceblue] (90:1.24) circle (0.08);
    \fill[targetgreen!80!black] (-30:1.24) circle (0.08);
    \fill[quotientorange!85!black] (-120:1.24) circle (0.08);

    \node[font=\scriptsize, align=center, text=sourceblue!90!black] at (0,0.05)
      {Major\\$0$--$4$--$7$};
  \end{scope}
\end{tikzpicture}
\caption{Left: Vibrating string harmonics showing why the octave ($2:1$, $+12$ semitones) and fifth ($3:2$, $+7$ semitones) align cleanly. Right: Wrapping the twelve semitones around the modular clock $\Z_{12}$, where $C$ major forms the rigid triangle $(0, 4, 7)$ --- rotating that same triangle clockwise builds every major chord.}
\label{fig:ratios-and-clock}
\end{figure}

\begin{intuitionbox}
Look at the blue triangle $(0, 4, 7)$ inside the twelve-hour clock in Figure~\ref{fig:ratios-and-clock}. Turn that rigid cardboard triangle two hours clockwise (add $+2$ to every corner): $(0, 4, 7)$ lands on $(2, 6, 9)$, which is $D$--$F\sharp$--$A$, a $D$ major chord. Because equal temperament spaces every hour on the clock by the same ratio $2^{1/12}$, turning a shape on the clock never stretches or warps its harmony.
\end{intuitionbox}

\section*{Exercises}

\begin{exercise}
Without looking at Figure~\ref{fig:two-octave-blueprint}, name the white key that sits (a) directly between the two chopsticks, (b) immediately to the right of the three-pronged fork, and (c) immediately to the left of the chopsticks. Check each answer on a keyboard or on Figure~\ref{fig:two-octave-blueprint}.
\end{exercise}

\begin{exercise}
Starting on the key $B$, count (a) one half step to the right, and (b) one whole step to the right. Which physical keys (white or black, and what letter names) do you land on? Now repeat both steps starting on $F$.
\end{exercise}

\begin{exercise}
Give two different directional names (enharmonic equivalents) for (a) the middle black key of the three-black-key fork, and (b) the white key immediately to the left of $C$.
\end{exercise}

\begin{exercise}
On the twelve-note clock $\Z_{12}$ ($0 = C, 1 = C\sharp, \dots, 11 = B$), rotate the $C$ major triangle $(0, 4, 7)$ clockwise by $+7$ hours (calculating each corner $\bmod 12$). Which three note names do you land on, and what major chord did you just build?
\end{exercise}
